\section{Compute 为什么尚未成为 Fungible Commodity}

在经济学和市场设计中，\term{fungibility} 指同一类别的一单位商品可以被另一单位替代，交易双方不必关心它来自哪个具体供应者。电网中的标准化电能接近这种体验，现代 AI compute 则高度异构：芯片架构、显存容量、互连、集群规模、软件版本、故障率、地域和合同条件都会影响任务能否运行。把所有资源折成 ``GPU-hour''，会隐藏大量不可替代性。

\subsection{一块 GPU 不是一度电}

训练任务需要一组同时满足的资源，而不是孤立芯片。可以把一份算力供给简化为资源向量
\[
\mathbf{c}=(n,\,m,\,b,\,l,\,s,\,a),
\]
其中 $n$ 表示可同时使用的加速器数量，$m$ 表示单卡可用显存，$b$ 表示互连带宽，$l$ 表示通信延迟，$s$ 表示软件与驱动兼容性，$a$ 表示可用性和故障恢复能力。只有当任务需求向量被整体满足时，这份供给才有实际价值；某一维度很高不能补偿另一维度完全不合格。

\begin{figure}[H]
\centering
\includegraphics[width=0.94\textwidth]{images/00-57-34-compute-nonfungible.jpg}
\caption{不同 GPU 和集群在价格、容量与可用特征上差异明显，compute 尚不能像标准电力一样互换。}
\figsource{Stanford Online 官方视频 00:57:34，`images/00-57-34-compute-nonfungible.jpg`。}
\end{figure}

\begin{knowledgebox}{读图：先问任务约束，再比较价格}
小模型推理可能在多种 GPU 上运行，大模型训练却要求大量同代卡、稳定高速互连、足够显存和长时间连续可用。图上的价格点若没有绑定这些属性，就无法说明哪种供给更便宜。对买方而言，不可替代性会抬高迁移成本；对卖方而言，异构规格又让统一市场、清算和容量规划更困难。
\end{knowledgebox}

\begin{warningbox}{首次使用提醒：HBM 与普通主存不是同一资源}
HBM 是 High Bandwidth Memory，高带宽存储器，通常与加速器通过先进封装紧密连接，为矩阵计算提供高带宽显存。服务器 DRAM 则主要作为 CPU 主存，容量、带宽和访问路径不同。即使两台机器都写着相同“内存总量”，训练中可用显存和通信拓扑仍可能完全不同。
\end{warningbox}

\subsection{Commodity 的五个工程条件}

讲者把可替代商品拆成五个条件：共同计量单位、标准交付接口、互连与池化、计量控制与结算、供应商之间可替换。这些条件不是经济学抽象名词，它们分别对应协议、调度器、资源描述、可靠性承诺和市场规则。Compute 若要从稀缺专用资产变成广泛可访问基础设施，五项能力需要一起成熟。

\begin{figure}[H]
\centering
\includegraphics[width=0.9\textwidth]{images/01-02-00-fungibility.jpg}
\caption{使一种资源成为 fungible commodity 的五项条件。}
\figsource{Stanford Online 官方视频 01:02:00，`images/01-02-00-fungibility.jpg`。}
\end{figure}

\begin{importantbox}{读图：五项条件对应什么系统能力}
共同单位要求资源规格可比较；标准交付接口要求任务能跨供应商提交；interconnection/pooling 要求分散容量能被组合；metering、control 与 settlement 要求准确记录使用、执行配额并完成结算；供应商可替换要求迁移不会重写全部软件。任何一项缺失，都可能让市场看起来有大量 GPU，实际却无法把容量交给需要它的任务。
\end{importantbox}

\begin{table}[H]
\centering
\begin{tabular}{L{0.2\textwidth}L{0.3\textwidth}L{0.36\textwidth}}
\toprule
条件 & 计算基础设施中的实现 & 当前难点 \\
\midrule
Common unit & 标准化性能、显存、互连和可用性描述 & 单一 FLOPs 或 GPU-hour 隐藏任务差异 \\
Delivery interface & 可移植 job spec、镜像、身份与数据接口 & 驱动、编译器、网络和权限碎片化 \\
Interconnection / pooling & 跨节点调度、容量池与故障迁移 & 延迟、地域、数据合规和拓扑约束 \\
Metering / control / settlement & 可信计量、配额、SLA、账单与清算 & 多租户隔离、性能抖动和责任归属 \\
Supplier substitution & 工作负载可跨云、跨芯片迁移 & 专有内核、服务 API 和数据重力 \\
\bottomrule
\end{tabular}
\caption{术语表：把 fungibility 转成可实现、可验证的系统接口。}
\end{table}

\subsection{标准为什么必须和制度一起出现}

技术标准规定资源怎样表示和连接，institutions 则负责协调采用、解决争议、执行承诺和保护公共利益。这里的 institution 可以是标准组织、行业联盟、市场运营者、监管机构或公共采购体系。只有文档没有执行，优势供应者可能拒绝互操作；只有强制命令没有可行接口，又会让规则停留在纸面。

\begin{figure}[H]
\centering
\includegraphics[width=0.94\textwidth]{images/01-03-00-standards-institutions.jpg}
\caption{历史基础设施通常由技术标准与协调制度共同推动可及性和稳定性。}
\figsource{Stanford Online 官方视频 01:03:00，`images/01-03-00-standards-institutions.jpg`。}
\end{figure}

\begin{knowledgebox}{读图：标准解决接口，制度解决协调}
AC/DC、电网调度、TCP/IP、电话互联或航空规则都不只是技术格式。它们需要测试、认证、争议处理、投资回收和公共责任安排。Compute 标准化同样会触及谁公开性能、谁承担中断、怎样跨境移动数据、如何避免市场操纵。讲者的核心主张是，fungibility 不能只靠更好的调度软件，还需要能够约束参与者并维护共同规则的机构。
\end{knowledgebox}

\begin{warningbox}{公共池化不等于忽略产权与安全}
把闲置容量纳入更大资源池，需要明确授权、隔离、计量、赔偿和退出机制。未经同意的强制共享会破坏投资激励与数据安全；完全不协调则可能让稀缺资源长期闲置或被少数主体囤积。制度设计的任务是把公共收益、供应可靠性和参与者权利写进可执行规则。
\end{warningbox}

\subsection{本章小结}

Compute 的不可替代性来自多维硬件、拓扑、软件和合同约束。Commodity 化需要共同单位、标准接口、池化、计量结算和供应商替换能力；这些技术能力又需要制度来推动采用、执行承诺和处理公共利益。至此，本讲从模型闭环走到了基础设施治理，也解释了课程为什么会同时邀请芯片、云、模型、应用与政策层的讲者。

